\documentclass[11pt]{article}

\usepackage[utf8]{inputenc}
\usepackage[T1]{fontenc}
\usepackage{geometry}
\usepackage{amsmath,amssymb,amsthm}
\usepackage{booktabs}
\usepackage{microtype}
\usepackage[hidelinks]{hyperref}

\geometry{margin=1in}

\theoremstyle{definition}
\newtheorem{definition}{Definition}
\newtheorem{game}{Game}
\theoremstyle{plain}
\newtheorem{proposition}{Proposition}
\newtheorem{corollary}[proposition]{Corollary}
\newtheorem{openq}{Open problem}

\newcommand{\Adv}{\mathsf{Adv}}
\newcommand{\A}{\mathcal{A}}
\newcommand{\B}{\mathcal{B}}
\newcommand{\Lrn}{\mathcal{L}}
\newcommand{\X}{\mathcal{X}}
\newcommand{\Y}{\mathcal{Y}}
\newcommand{\Q}{\mathcal{Q}}
\newcommand{\M}{\mathcal{M}}
\newcommand{\Oracle}{\mathcal{O}}
\newcommand{\INFER}{\textsc{Infer}}
\newcommand{\LEARN}{\textsc{Learn}}
\newcommand{\GEN}{\textsc{GenerateData}}

\title{Extraction games with a priced adversary\\
\large An intermediate report for the survey's formalization}
\author{}
\date{2026-08-22, second pass}

\begin{document}
\maketitle

\section{What this report is}

This report turns a set of proposed games into definitions, states what each one
settles and what it does not, and identifies the reduction targets available to
them. It is written to be reused. Sections~\ref{sec:g2} to \ref{sec:price} are
candidate text for the survey's formalization section, and
Section~\ref{sec:reduce} is the part that must be settled before any of it carries
a theorem.

Three commitments hold throughout. Every game names a baseline, because a success
probability without one is not a measurement. Every game prices both parties,
because the object of interest is a ratio. And every claim is marked as proved
here, assumed, or open.

\paragraph{Scope.} The games below assume a metric output space in which a
tolerance $\varepsilon$ is meaningful, so they describe regression and
real-valued-output models. Under a $0/1$ output metric, Propositions~\ref{prop:collapse}
and~\ref{prop:sep} are vacuous, because no useful Lipschitz constant exists. The
classification case needs a different treatment and this report does not give one.

\section{Notation}
\label{sec:notation}

Fix an input space $\X$ with metric $d_\X$ and an output space $\Y$ with metric
$d_\Y$. We use the pipeline decomposition of Villa et al.~\cite{VDKMR26}, in which
a system is a tuple $(\GEN, \LEARN, \INFER)$ carrying a source to a corpus, a
corpus to a model, and a model with an input to a result. We write $\mathsf{ATK}$
for their capability set, with elements $\mathsf{see\_model}$,
$\mathsf{black\_box}$, $\mathsf{inject\_learn}$ and $\mathsf{inject\_infer}$.

Everything is indexed by a security parameter $n$, and this is not a formality.
Without it the hardness statement in Section~\ref{sec:price} cannot be written.
Let $\M_n$ be a distribution over models $f : \X_n \to \Y_n$, induced by running
$(\GEN, \LEARN)$ on a source. The learner $\Lrn$ holds $f \sim \M_n$. We write
$C_{\INFER}$ for the cost of one honest evaluation of $f$, and measure every
party's budget as a dimensionless multiple of it.

\section{The extrapolation game with an adversary-chosen probe}
\label{sec:g1}

The first proposal is the most natural. The adversary queries, then names a point
it did not query and predicts the model's value there.

\begin{game}[$\mathsf{Extr}^{\mathrm{adv}}$]
\label{game:g1}
Parameters: a query budget $m$, a probe radius $\Delta$, a tolerance $\varepsilon$.
\begin{enumerate}
  \item $\A$ chooses $Q = \{q_1,\dots,q_m\} \subset \X$ and receives
        $R = \{r_i = f(q_i)\}_{i \le m}$.
  \item $\A$ outputs $(q', r')$ with $q' \notin Q$ and $d_\X(q', Q) \le \Delta$.
  \item $\A$ wins if $d_\Y(r', f(q')) \le \varepsilon$.
\end{enumerate}
\end{game}

This game is not usable as stated, and the reason constrains every later version.

\begin{proposition}[The probe collapses]
\label{prop:collapse}
Let $f$ be $L$-Lipschitz from $(\X,d_\X)$ to $(\Y,d_\Y)$. Suppose the rules of
Game~\ref{game:g1} admit some $q'$ with $0 < d_\X(q', q_i) \le \delta$ for a
queried $q_i$, and that $L\delta \le \varepsilon$. Then the specific adversary that
queries one point $q_i$, names such a $q'$ and outputs $r' = r_i$ wins with
probability $1$, using one query and no computation beyond naming $q'$.
\end{proposition}

\begin{proof}
$d_\Y(r_i, f(q')) = d_\Y(f(q_i), f(q')) \le L\, d_\X(q_i, q') \le L\delta \le
\varepsilon$.
\end{proof}

Proposition~\ref{prop:collapse} is an existence statement about one adversary, and
that is enough to reject the game. Where such a $q'$ is admissible the game
measures the modulus of continuity of $f$ rather than the adversary. Two repairs
exist. The first forces the probe away from the queried set, requiring
$d_\X(q',Q) \ge \delta_{\min}$ with $L\delta_{\min} > \varepsilon$, which makes the
difficulty depend on $L$, a property of the model that the adversary may not know
and that the isoperimetry line studies in its own right~\cite{BS21}. The second
removes the adversary's choice, and is the one developed here.

\section{The learner-chosen probe}
\label{sec:g2}

Let the learner name the challenge points. The geometry disappears, because success
is measured where the learner chose rather than where the adversary argued was
close enough.

\begin{game}[$\mathsf{Extr}^{\mathrm{lrn}}$]
\label{game:g2}
Parameters: a query budget $m$, a challenge count $k$, a tolerance $\varepsilon$, a
challenge distribution $\Q$ over $k$-subsets of $\X$, a mode
$\mathsf{mode} \in \{\mathsf{commit},\mathsf{adapt}\}$. Sample $f \sim \M_n$.
\begin{enumerate}
  \item If $\mathsf{mode} = \mathsf{commit}$, $\Lrn$ samples $Q^* \sim \Q$ and
        publishes a commitment to it.
  \item $\A$ chooses $Q = \{q_1,\dots,q_m\}$ and receives $R = \{f(q_i)\}_{i\le m}$.
  \item If $\mathsf{mode} = \mathsf{adapt}$, $\Lrn$ samples $Q^* \sim \Q(Q)$ now.
        $\Lrn$ sends $Q^* = \{q^*_1,\dots,q^*_k\}$, opening the commitment in the
        committed mode.
  \item $\A$ outputs $(r^*_1,\dots,r^*_k)$.
  \item $\A$ wins if $d_\Y(r^*_j, f(q^*_j)) \le \varepsilon$ for every $j \le k$.
\end{enumerate}
\end{game}

The two modes differ in what they measure. In the committed mode the challenge is
independent of the adversary's queries, so the game measures generalization from a
fixed test set. In the adaptive mode the learner may place challenges where the
adversary did not look, so the game measures coverage. That the adaptive mode is
harder for the adversary is intuitive and is not proved here; a coupling argument
would settle it, and Open Problem~\ref{oq:adapt} records it.

\subsection{The advantage, and why the baseline is part of the definition}

The baseline must be defined so that it cannot smuggle in the model. If $f$ is
fixed and the baseline maximizes over all algorithms that ``do not access $f$'',
the maximizing algorithm can simply have $f$ hardwired, the baseline equals the
honest success probability, and the advantage is non-positive for every adversary.
That is the same degeneracy that makes the marginal-advantage form of membership
advantage a no-op. The repair is to draw the model from a distribution and measure
the baseline against that distribution, which is what Balle, Cherubin and Hayes do
by drawing the target record from a prior~\cite{BCH22}.

\begin{definition}[Extraction advantage]
\label{def:adv}
Let $\B$ range over algorithms that receive the public description of $\M_n$, the
task, and the challenge set $Q^*$, but neither $f$ nor the transcript $(Q,R)$. Then
\[
\Adv^{\mathsf{Extr}}_{\M_n,\Q}(\A)
= \Pr\big[\A \text{ wins}\big] - \max_{\B} \Pr\big[\B \text{ wins}\big],
\]
both probabilities over $f \sim \M_n$, over $\Q$, over the coins of $\A$ and $\B$,
and over any randomness in $\INFER$.
\end{definition}

Because $f$ is random and $\B$ is quantified before $f$ is drawn, no $\B$ can
hardwire $f$, and the baseline is the best model-free predictor for the
\emph{family} rather than for one model. Definition~\ref{def:adv} still differs
from reconstruction robustness in exactly one place, and the difference is the
point of this report: Balle, Cherubin and Hayes place no bound on the adversary's
computation, quantifying over any reconstruction map, while
Section~\ref{sec:price} adds one. Their $R$ is the attack, and the criterion is a
separate error function with tolerance $\eta$.

\begin{proposition}[The raw success probability is not a measurement]
\label{prop:baseline}
There exist a task, a model family and a challenge distribution for which a
model-free $\B$ attains success probability strictly greater than that of a
published attack with full query access.
\end{proposition}

\begin{proof}[Status and evidence]
Empirical, and established for membership inference rather than for extraction.
Das et al.\ exhibit it on eight published membership-inference benchmarks, where
classifiers with no model access outperform the best reported attack that queries
the model~\cite{DZT25}. Meeus et al.\ report the same effect and identify the
carrier, a distribution shift between the two halves of each benchmark
\cite{MSJFRM25}. Whether the same construction arises for extraction challenge
distributions is asserted here and not exhibited.
\end{proof}

\subsection{Separated challenges}

The extension to $k$ challenges is meant to measure utility rather than one lucky
point. Without a constraint on how the challenges sit relative to one another it
measures neither.

\begin{proposition}[Unseparated challenges collapse to one]
\label{prop:sep}
Let $f$ be $L$-Lipschitz. Suppose the support of $\Q$ is contained in a ball
$B(c,\rho)$ whose centre $c$ is known to the adversary before it queries, and
suppose $L\rho \le \varepsilon$. Then the adversary that queries $c$ once and
answers every challenge with $f(c)$ wins Game~\ref{game:g2} with probability $1$,
for every $k$, in either mode.
\end{proposition}

\begin{proof}
For each $j$, $d_\Y(f(c), f(q^*_j)) \le L\,d_\X(c, q^*_j) \le L\rho \le
\varepsilon$, and the win condition is a conjunction over $j$.
\end{proof}

The hypothesis that the ball is fixed and known before step 2 is necessary, because
the adversary queries before $Q^*$ is revealed. Two constraints repair the game, and
they are not equivalent.

\begin{definition}[$\Delta$-separated challenge distribution]
\label{def:packing}
$\Q$ is $\Delta$-separated if every $Q^*$ in its support satisfies
$\min_{j \ne j'} d_\X(q^*_j, q^*_{j'}) \ge \Delta$. The $\Delta$-packing number of
$\X$ upper bounds the challenge count $k$.
\end{definition}

\begin{definition}[$h$-unpredictable challenge distribution]
\label{def:minent}
$\Q$ is $h$-unpredictable if $H_\infty(Q^*) \ge h$, so no algorithm guesses $Q^*$
with probability better than $2^{-h}$.
\end{definition}

Definition~\ref{def:packing} is geometric and rules out
Proposition~\ref{prop:sep}. Definition~\ref{def:minent} is information-theoretic
and rules out an adversary that precomputes answers for a predictable challenge
set. A game that prices an adversary needs both, because packing alone permits a
fixed challenge set and min-entropy alone permits a clustered one. The pair is the
analogue of requiring distinct and unpredictable challenge messages in an
indistinguishability game.

\section{The decision variant}
\label{sec:ind}

The search game asks the adversary to produce a value. The crypto-native form asks
it to distinguish, giving an advantage against a coin rather than against a
maximization over baselines.

\begin{game}[$\mathsf{Extr}\text{-}\mathsf{IND}$]
\label{game:ind}
As Game~\ref{game:g2} through step 3. Then $\Lrn$ samples $b \leftarrow \{0,1\}$.
If $b = 0$ it returns $y_j = f(q^*_j)$ for each $j$. If $b = 1$ it returns
$y_j \sim \mathcal{R}(q^*_j)$, independently across $j$, for a reference
distribution $\mathcal{R}$ independent of $f$. $\A$ outputs $b'$ and wins if
$b' = b$. Its advantage is $\Adv^{\mathsf{IND}}(\A) = |\Pr[b'=b] - \tfrac12|$.
\end{game}

\begin{proposition}[Search implies decision]
\label{prop:s2d}
Let $\A$ win Game~\ref{game:g2} with probability $p$. Define
\[
q \;=\; \mathbb{E}_{f, Q^*, \text{coins of } \A}
\left[\; \prod_{j \le k}
\Pr_{y_j \sim \mathcal{R}(q^*_j)}\!\left[\, d_\Y(y_j, r^*_j) \le \varepsilon \,\right]
\right],
\]
the probability that a full reference vector falls within tolerance of $\A$'s own
predictions. Then there is a distinguisher $\A'$ with
$\Adv^{\mathsf{IND}}(\A') = \tfrac12|p - q|$, running in the time of $\A$ plus $k$
distance evaluations.
\end{proposition}

\begin{proof}
$\A'$ runs $\A$ to completion on $(Q^*, R)$ to obtain $(r^*_1,\dots,r^*_k)$, and
only then reads $(y_1,\dots,y_k)$. It outputs $b'=0$ if $d_\Y(y_j, r^*_j) \le
\varepsilon$ for all $j$, and $b'=1$ otherwise. Because $\A$ has finished before
$y$ is read, the vector $r^*$ is independent of $y$ given $Q^*$, and because the
$y_j$ are independent across $j$ the acceptance probability under $b=1$ is the
product above. Hence $\Pr[b'=0 \mid b=0] = p$ and $\Pr[b'=0 \mid b=1] = q$, so
$\Pr[b'=b] = \tfrac12(p + 1 - q)$ and the advantage is $\tfrac12|p-q|$.
\end{proof}

The converse is the interesting direction and is not available in general. The
nearest published statement is Alvarez-Melis, Garg and Kalai~\cite{AGK22}, and it
must be quoted carefully because the paper describes the same reduction from both
ends. Its introduction calls the result ``a polynomial-time reduction from
likelihood maximization to next-token distinguishability'', while the body
concludes that its Algorithm 1 with Theorem 1 and Lemma 3 ``yields a
polynomial-time reduction from distinguishing distributions to maximum likelihood
estimation for sequential models''. The unambiguous operational statement is that
Algorithm 1 converts a next-token distinguisher into a likelihood improvement, so a
distinguisher yields a better generator, and no converse is stated in that model.
Two further boundaries hold. The two-directional result is in a different section
and is not polynomial time, the paper noting that repeating its efficient step
gives an exponential-time algorithm overall. And no distinguisher anywhere in the
paper carries a runtime bound, although the next-token result does fix a
distinguisher class supplied by an oracle.

\begin{openq}
\label{oq:converse}
Is there a polynomial-time reduction from $\mathsf{Extr}\text{-}\mathsf{IND}$ to
$\mathsf{Extr}^{\mathrm{lrn}}$ under a $\Delta$-separated, $h$-unpredictable
challenge distribution? A positive answer makes the decision game the definition of
record and the search game a corollary.
\end{openq}

\begin{openq}
\label{oq:adapt}
Prove by coupling that the adaptive mode of Game~\ref{game:g2} is at least as hard
for the adversary as the committed mode, or exhibit a $\Q$ where it is not.
\end{openq}

\section{Pricing the two parties}
\label{sec:price}

Every game above is a predicate on a transcript. It becomes a security notion only
when both parties carry a budget, and the budget must be written so that the
adversary cannot inflate its own allowance.

Measure cost in evaluations of $f$, so all quantities are dimensionless. The
adversary's total cost is $t_\A(n) = m + T_\A/C_{\INFER}$, where $m$ is its query
count and $T_\A$ its own computation. The honest reference cost is $k(n)$, the
number of evaluations the learner needs to produce the $k$ challenge answers it is
adjudicating against.

The reference cost deliberately excludes both the $m$ queries the learner answers
and the cost of training $f$. Excluding the queries matters, because a threshold
that depended on $m$ would rise as the adversary queried more, letting the
adversary raise its own admission bar. Excluding training matters because a
threshold dominated by pre-training is met by no query-bounded adversary and the
definition would be vacuous.

\begin{definition}[$\delta$-extraction hardness]
\label{def:fg}
A family $\{(\M_n, \Q_n, k(n), \varepsilon(n))\}_n$ is $\delta$-extraction hard if
for every adversary $\A$ with $t_\A(n) < k(n)^{1+\delta}$ for all large $n$,
\[
\Adv^{\mathsf{Extr}}_{\M_n,\Q_n}(\A) \;\le\; \mathrm{negl}(n).
\]
\end{definition}

The name carries only $\delta$, deliberately. Calling it
$(\varepsilon,\delta)$-hardness would collide with Brzuska and
Couteau~\cite{BC22}, where $\varepsilon$ bounds the inverter's success
probability, whereas $\varepsilon$ here is the game's tolerance. Three cautions
apply and the survey must state all three.

\begin{enumerate}
  \item \textbf{Three conventions, not one.} Brzuska and Couteau use non-uniform
        circuit size with the honest cost raised to $1+\delta$~\cite{BC22}.
        LaVigne, Lincoln and Vassilevska Williams use uniform word-RAM time and
        name the gap as a constant $\delta$ separating honest $T^{1-\delta}$ from
        the threshold $T$~\cite{LLW19}. Ball, Rosen, Sabin and
        Vasudevan~\cite{BRSV17} index their error by the adversary's own discount
        as well as by $n$, so one number does not describe their notion. Quoting
        one convention's $\delta$ in another's units is an error, and
        Definition~\ref{def:fg} is in the first convention.
  \item \textbf{The known ceiling is a query bound.} Barak and
        Mahmoody~\cite{BM09} show an adversary making $O((n/\delta)^2)$
        \emph{oracle queries} breaks any $n$-query key exchange in the random
        oracle model, and note their adversary runs efficiently in a relativized
        world where $\mathrm{P} = \mathrm{NP}$. That forbids a super-quadratic
        \emph{query} gap in an idealized model. It says nothing about time gaps or
        structured assumptions, so Definition~\ref{def:fg} is not capped at
        $\delta = 1$ by it.
  \item \textbf{The exponent is on a dimensionless count.} Raising a cost in
        seconds or FLOPs to the power $1+\delta$ is not invariant under a change of
        unit. Measuring in evaluations of $f$ removes the ambiguity, and matches
        Brzuska and Couteau, whose $|f|$ is a circuit size.
\end{enumerate}

\section{Model stealing under an abstracted distance}
\label{sec:steal}

Extraction of the function itself is the same game with the output replaced by a
model. The difficulty is that the natural distance is the wrong one. Two networks
that agree everywhere differ in weights, and a transformer's trunk is largely a
product of pre-training that no attacker needs to steal.

Let $\rho$ be a \emph{read-out map} carrying a model to the part of it that
matters, and let $D_\rho(f,f') = d(\rho(f), \rho(f'))$.

\begin{game}[$\mathsf{Steal}_\rho$]
\label{game:steal}
As Game~\ref{game:g2} through step 2. Then $\A$ outputs a model $f'$, and wins if
$D_\rho(f, f') \le \varepsilon$.
\end{game}

\begin{proposition}[Stealing implies answering]
\label{prop:steal}
Let $\rho$ restrict a model to the support of $\Q$ and let $d$ be the supremum
distance. If $\A$ wins $\mathsf{Steal}_\rho$ with tolerance $\varepsilon$, then
evaluating $f'$ at each challenge wins Game~\ref{game:g2} with the same
$\varepsilon$, at the cost of $k$ evaluations of $f'$.
\end{proposition}

\begin{proof}
$d_\Y(f'(q^*_j), f(q^*_j)) \le \sup_{q \in \mathrm{supp}\,\Q} d_\Y(f'(q),f(q)) =
D_\rho(f,f') \le \varepsilon$ for every $j$.
\end{proof}

The converse fails without a covering condition. Answering $k$ separated challenges
pins $f$ at $k$ points, and $k$ points determine $f$ on the support only if they
form an $\alpha$-net for a modulus of continuity bounding $f$. This is where
Definition~\ref{def:packing} does work beyond ruling out
Proposition~\ref{prop:sep}: the packing number bounds how many separated challenges
exist, and the covering number bounds what $k$ challenges can certify.

\begin{openq}
For which $(\rho, \Q, k)$ does winning Game~\ref{game:g2} on a $\Delta$-separated
challenge set imply winning $\mathsf{Steal}_\rho$? A quantitative answer converts a
query-complexity bound into a model-extraction bound and back.
\end{openq}

The choice of $\rho$ decides what the theorem is about and must be stated rather
than assumed. Jagielski et al.~\cite{JCBKP20} separate the two natural choices as
accuracy and fidelity objectives, and show they conflict for an imperfect victim,
since a high-fidelity copy must reproduce the victim's errors and a high-accuracy
copy must not. They also dismiss exact weight recovery as the wrong target, because
ReLU networks sit in large behavioural equivalence classes closed under neuron
rescaling, permutation and dead-neuron insertion. Any $\rho$ used here must be
invariant under those symmetries or it measures an artifact.

\section{The challenge generator as an oracle}
\label{sec:oracle}

In the adaptive mode the learner chooses $Q^*$ after seeing the adversary's
queries. That choice is an algorithm, and leaving it unspecified hides a party
whose power decides the game.

Model it as an auxiliary oracle $\Oracle_{\mathrm{chal}}$ in the sense of Villa et
al.~\cite{VDKMR26}, with its own $(\GEN, \LEARN, \INFER)$ and capability set. It
need not be a language model. A derivative of $f$, or a smaller model trained on
the same task, suffices.

Three consequences follow.

\begin{enumerate}
  \item The capability set becomes the access coordinate of the survey's frame
        directly. $\mathsf{black\_box}$ and $\mathsf{see\_model}$ separate the query
        and white-box variants of every game above, and $\mathsf{inject\_learn}$
        opens the poisoning variant in which $\A$ influences $f$ before the game.
  \item Villa et al.\ separate completeness, measured with no adversary, from
        security, measured against one. Game~\ref{game:g2} needs the same split,
        because a challenge set the honest learner cannot itself answer within
        $\varepsilon$ makes the game unwinnable for reasons unrelated to the
        adversary.
  \item A filtering impossibility bears on how cheap $\Oracle_{\mathrm{chal}}$ can
        be, and its attribution matters. The result is Ball et al.~\cite{BGGKRR25},
        which Villa et al.\ discuss rather than prove. It is \emph{existential}:
        the abstract states that there exist language models for which no efficient
        prompt filter works, and the construction assumes the model shares a secret
        key with the attacker. It does not say that every judge must match the
        complexity of its predicate. Read at its actual strength it suggests, and
        does not establish, that a challenge generator too weak to compute the
        property it tests cannot place hard challenges, which would put
        $\Oracle_{\mathrm{chal}}$'s own cost inside the honest budget.
\end{enumerate}

\section{Where to reduce}
\label{sec:reduce}

The proposal asks what to reduce to when there is no instance in hand. Four answers
exist.

\paragraph{A definitional target with no reduction.} A game may stand as a
definition before any instantiation. What this buys is a statement other results
can be mapped into, which is the survey's goal. What it costs is that no theorem
follows. The programme's own standing constraint is that the machinery here is
off-the-shelf, so the contribution must be a theorem at a stated $\delta$ and not a
definition. This option therefore satisfies the survey and not the thesis, and
saying so is part of the report.

\paragraph{Learning-theoretic exponents.} Two distinct papers are involved and
conflating them is easy. Damian, Lee and Bruna~\cite{DLB25} give
$n = \Theta(d^{1 \vee k^*/2})$ for Gaussian multi-index models with the generative
leap exponent, with a low-degree lower bound and a matching efficient upper bound
whose constant $C(P,K)$ is conjectured exponentially large in the index dimension.
The detection caveat belongs to the single-index paper, Damian, Pillaud-Vivien, Lee
and Bruna~\cite{DPLB24}, whose Remark 3.3 says its theorem primarily measures a
detection threshold, and whose low-degree half is conditional on its Conjecture
3.4. An adversary that must recover a value is doing estimation, whose threshold
those authors leave open. Separately, multi-pass gradient descent with batch reuse
escapes the information and leap exponents~\cite{DTAPZK24}, so these price
single-pass training rather than intrinsic difficulty.

\paragraph{Planted problems.} Planted clique, planted subgraph and continuous LWE
supply average-case hardness with a long record. The caution is that restricted-class
evidence and the security statement are different objects. Bogdanov, Jones, Rosen
and Zadik~\cite{BJRZ24} build private simultaneous messages and secret sharing whose
security quantifies over size-bounded circuits, with low-degree analysis offered as
evidence for the underlying conjectures. No security definition there quantifies
over a low-degree adversary. Buhai, Hsieh, Jain and Kothari~\cite{BHJK25} disprove
the quasi-polynomial low-degree conjecture in the regime it was stated, so
low-degree-bounded and quasi-polynomial-bounded are provably different. Two clauses
travel with that: their polynomial-time distinguisher is for the rectangular
generalisation, while in the square symmetric regime it is
$n^{O(\log^{1/(k-1)} n)}$, and the positive counterpart~\cite{HKKLMT26} must be
cited alongside.

\paragraph{Fine-grained hard problems.} Average-case $k$-SUM and Zero-$k$-Clique
carry fine-grained hardness and are the established route to a statement shaped
like Definition~\ref{def:fg}~\cite{LLW19,BRSV17}. This is the target that matches
the quantity, because the conclusion wanted here is a ratio and these are the only
assumptions stated as ratios. The cost is a theorem about a constructed family
rather than a deployed model, a boundary the survey already commits to stating.

\section{What is already occupied}
\label{sec:prior}

\begin{itemize}
  \item \textbf{The criterion, tolerance and baseline exist.} Balle, Cherubin and
        Hayes~\cite{BCH22} define reconstruction robustness with an error function,
        a tolerance $\eta$, and a baseline given by the best attack that ignores the
        model. What is missing there is any bound on the adversary's computation.
  \item \textbf{An extraction adversary priced by algorithm class exists, in a
        security paper.} Jagielski et al.~\cite{JCBKP20} state as their Theorem 3,
        imported from Das et al., that model extraction approaches expressible as
        statistical-query algorithms require $\exp(O(h))$ samples for fidelity
        extraction of depth-$h$ random networks, and note that gradient descent and
        PCA are such algorithms. This is closer to the present proposal than
        anything else found, and any claim of novelty must be stated against it.
  \item \textbf{Charging the adversary compute per query is deployed.} Dziedzic et
        al.~\cite{DKLP22} wrap every prediction in a proof-of-work puzzle whose
        difficulty is calibrated to how much information the client has already
        extracted, reporting under $2\times$ slowdown for legitimate users and up to
        $100\times$ for attackers. It carries no definition and no theorem, which is
        the gap, but ``prices the attack in compute'' without that qualifier is
        already published.
  \item \textbf{The query budget is already a definitional parameter.}
        Chandrasekaran et al.~\cite{CCGJY19} give the pre-2023 literature its
        closest thing to a formal extraction definition, with the query budget
        inside it. Their reduction turns a query-synthesis active learner into an
        extraction adversary, so extraction reduces to active learning and the
        complexity figures travel in the upper-bound direction only. The converse,
        that a class hard to actively learn is hard to extract, is neither proved
        nor claimed.
  \item \textbf{Extraction as cryptanalysis exists.} Carlini, Jagielski and
        Mironov~\cite{CJM20} recover ReLU parameters by differential cryptanalysis
        from real-valued oracle access. The attack uses polynomially many queries
        but work exponential in layer width for sign recovery, and it is empirical.
  \item \textbf{The success definition is reported missing from inside.} Shafran et
        al.~\cite{SSEP24} state that they are aware of no model-extraction
        literature that formally defines what makes an attack successful. The scope
        of that statement is accuracy-driven extraction, with fidelity and
        functionally equivalent extraction excluded, and the work is a preprint.
  \item \textbf{The attacker-versus-honest ratio is stated, and parameterized by a
        constant but never by an exponent.} Proof-of-learning requires that forging
        cost at least as much as honest training, in its Property 2 and its
        Definition 2, offering a conjecture in place of a proof~\cite{JYCDTCP21}.
        Its Property 1 is a different inequality, bounding verification against
        training. Later work does index the ratio, by a multiplicative constant in
        Fang et al.~\cite{ZFWCSXZ26} and by efficiency and hardness parameters in
        the compute-monitoring proposal of Shavit~\cite{Sha23}. What no work does is
        index it by an exponent, which is the only form Definition~\ref{def:fg}
        adds.
\end{itemize}

\section{Open problems, in the order they block progress}

\begin{enumerate}
  \item Open Problem~\ref{oq:converse}, the polynomial-time converse of
        Proposition~\ref{prop:s2d}.
  \item The covering condition under which Game~\ref{game:g2} implies
        $\mathsf{Steal}_\rho$, quantitatively in $k$ and $\Delta$.
  \item A read-out map $\rho$ invariant under the ReLU symmetries that Jagielski et
        al.\ name, so that $D_\rho$ measures function and not parameterization.
  \item An instance meeting Definition~\ref{def:fg} at an explicit $\delta$,
        reduced to average-case $k$-SUM or Zero-$k$-Clique.
  \item A separation between Definition~\ref{def:fg} and the statistical-query
        bound of Jagielski et al.\ Theorem 3, exhibiting a family hard for a
        compute-bounded adversary and easy for an unbounded statistical-query one,
        or the reverse. Without it the exponent form adds no theorem.
  \item The classification case, where the tolerance $\varepsilon$ and the
        Lipschitz arguments of Section~\ref{sec:g1} do not apply.
\end{enumerate}

\bibliographystyle{plain}
\bibliography{../refs}

\end{document}
